\documentclass[10pt,a4paper]{amsart}
\usepackage{amsmath,amssymb,amsthm,graphicx}
% Retained preamble for the Stacks Project chapter snapshot in this directory.
% The chapter's own preamble.tex is a copy of the shared Stacks Project preamble; the block
% below is sliced MECHANICALLY, line for line, from it. Nothing was added, removed or
% redefined; the wrapper class, the xr-hyper/hyperref wiring and the externaldocument list
% are omitted because the assembler supplies the document class.
%
% lines taken from preamble.tex: 3, 5, 9, 11, 12, 156, 157, 158, ... (51 lines total)
%
% package coverage restored from the chapter's own preamble (2026-10-13): xy, xy-2cell

\usepackage[all]{xy}
\xyoption{2cell}
\UseAllTwocells

\newenvironment{reference}{\comment}{\endcomment}
\newenvironment{slogan}{\comment}{\endcomment}
\newenvironment{history}{\comment}{\endcomment}
\theoremstyle{plain}
\newtheorem{theorem}[subsection]{Theorem}
\newtheorem{proposition}[subsection]{Proposition}
\newtheorem{lemma}[subsection]{Lemma}
\theoremstyle{definition}
\newtheorem{definition}[subsection]{Definition}
\newtheorem{example}[subsection]{Example}
\newtheorem{exercise}[subsection]{Exercise}
\newtheorem{situation}[subsection]{Situation}
\theoremstyle{remark}
\newtheorem{remark}[subsection]{Remark}
\newtheorem{remarks}[subsection]{Remarks}
\numberwithin{equation}{subsection}
\def\lim{\mathop{\mathrm{lim}}\nolimits}
\def\colim{\mathop{\mathrm{colim}}\nolimits}
\def\Spec{\mathop{\mathrm{Spec}}}
\def\Hom{\mathop{\mathrm{Hom}}\nolimits}
\def\Ext{\mathop{\mathrm{Ext}}\nolimits}
\def\SheafHom{\mathop{\mathcal{H}\!\mathit{om}}\nolimits}
\def\SheafExt{\mathop{\mathcal{E}\!\mathit{xt}}\nolimits}
\def\Sch{\mathit{Sch}}
\def\Mor{\mathop{\mathrm{Mor}}\nolimits}
\def\Ob{\mathop{\mathrm{Ob}}\nolimits}
\def\Sh{\mathop{\mathit{Sh}}\nolimits}
\def\NL{\mathop{N\!L}\nolimits}
\def\CH{\mathop{\mathrm{CH}}\nolimits}
\def\proetale{{pro\text{-}\acute{e}tale}}
\def\etale{{\acute{e}tale}}
\def\QCoh{\mathit{QCoh}}
\def\Ker{\mathop{\mathrm{Ker}}}
\def\Im{\mathop{\mathrm{Im}}}
\def\Coker{\mathop{\mathrm{Coker}}}
\def\Coim{\mathop{\mathrm{Coim}}}
\DeclareMathSymbol{\boxtimes}{\mathbin}{AMSa}{"02}
\def\QCohstack{\mathcal{QC}\!\mathit{oh}}
\def\Cohstack{\mathcal{C}\!\mathit{oh}}
\def\Spacesstack{\mathcal{S}\!\mathit{paces}}
\def\Quotfunctor{\mathrm{Quot}}
\def\Hilbfunctor{\mathrm{Hilb}}
\def\Curvesstack{\mathcal{C}\!\mathit{urves}}
\def\Polarizedstack{\mathcal{P}\!\mathit{olarized}}
\def\Complexesstack{\mathcal{C}\!\mathit{omplexes}}
\def\Pic{\mathop{\mathrm{Pic}}\nolimits}
\def\Picardstack{\mathcal{P}\!\mathit{ic}}
\def\Picardfunctor{\mathrm{Pic}}
\def\Deformationcategory{\mathcal{D}\!\mathit{ef}}
\usepackage{cleveref}
\crefname{originallemma}{original lemma}{original lemmas}
\Crefname{originallemma}{Original lemma}{Original lemmas}
\makeatletter
\newlabel{lemma-acyclicity}{{01GD}{source}{}{sourcepointer-1}{}}
\newlabel{lemma-acyclicity@cref}{{[originallemma][0][]01GD}{[1][0][]source}{}{}{}}
\makeatother
\crefname{originallemma}{original lemma}{original lemmas}
\Crefname{originallemma}{Original lemma}{Original lemmas}
\makeatletter
\newlabel{lemma-compare-cosk0}{{01GC}{source}{}{sourcepointer-2}{}}
\newlabel{lemma-compare-cosk0@cref}{{[originallemma][0][]01GC}{[1][0][]source}{}{}{}}
\makeatother
\crefname{originallemma}{original lemma}{original lemmas}
\Crefname{originallemma}{Original lemma}{Original lemmas}
\makeatletter
\newlabel{simplicial-example-fibre-products-simplicial-object}{{016E}{source}{}{sourcepointer-3}{}}
\newlabel{simplicial-example-fibre-products-simplicial-object@cref}{{[originallemma][0][]016E}{[1][0][]source}{}{}{}}
\makeatother
\crefname{originallemma}{original lemma}{original lemmas}
\Crefname{originallemma}{Original lemma}{Original lemmas}
\makeatletter
\newlabel{simplicial-lemma-cosk-minus-one}{{018I}{source}{}{sourcepointer-4}{}}
\newlabel{simplicial-lemma-cosk-minus-one@cref}{{[originallemma][0][]018I}{[1][0][]source}{}{}{}}
\makeatother
\crefname{originallemma}{original lemma}{original lemmas}
\Crefname{originallemma}{Original lemma}{Original lemmas}
\makeatletter
\newlabel{simplicial-remark-augmentation}{{018J}{source}{}{sourcepointer-5}{}}
\newlabel{simplicial-remark-augmentation@cref}{{[originallemma][0][]018J}{[1][0][]source}{}{}{}}
\makeatother
\title{Acyclicity of simplicial presheaves satisfying the hypercovering conditions}
\author{The Stacks Project Authors}
\date{stacks-project:hypercovering.tex:local-snapshot-20261012; GNU FDL 1.2 or later}
\begin{document}
\maketitle
\noindent\textit{Editorial note.} Statement span is lines 934-960 of sources/stacks-hypercovering/pinned-original-source.tex; proof span is lines 962-998, the complete author proof (35 body lines).\par\medskip
\noindent\textit{Editorial note.} Navigation stubs: 5 cited-but-not-included labels are supplied as non-load-bearing pointers with their pinned tags where the snapshot tag index has them.\par\medskip
\noindent\textit{Editorial note.} extra\_packages is empty on purpose: the retained preamble already loads xy with options and a second declaration clashes in the assembler wrapper.\par\medskip
\noindent\textit{Editorial note.} Text is retained verbatim, including the author's own typos and placeholders.\par\medskip
\noindent\textit{Editorial note.} Queue index 12 of queue\_hypercovering.json (pb=1521, stmt 934-960, proof 962-998).\par\medskip
\noindent\textit{Editorial note.} stacks-hypercovering wave-1 (hc01): fresh adjudication accept (2026-10-14); pool per standing fallback order (hypercovering -> pic -> spaces-more-cohomology) after md01 closed the moduli pool; hypercovering.tex 53 labels/34 with proof judged 26 accept / 8 reject / 14 pb=0 env-dup rows excluded; spans cite-free and reference/slogan/history-free so no patchenv needed.\par\medskip
\noindent\textit{Allowed standard prerequisites (editorial).}\begin{itemize}
\item Explicit prerequisite boundary. The cited inputs are prerequisites; the argument in the retained span is not replaced by any cited result.
\end{itemize}\medskip
\ifdefined\hypertarget\hypertarget{sourcepointer-1}{}\fi
\noindent\textit{Source reference.} Original lemma 01GD: sources/stacks-hypercovering/pinned-original-source.tex, source lines 880--891. Cited by the retained statement or proof but not part of the collected spans; supplied only as a navigation pointer. The argument does not depend on its content.\par\medskip
\ifdefined\hypertarget\hypertarget{sourcepointer-2}{}\fi
\noindent\textit{Source reference.} Original lemma 01GC: sources/stacks-hypercovering/pinned-original-source.tex, source lines 808--824. Cited by the retained statement or proof but not part of the collected spans; supplied only as a navigation pointer. The argument does not depend on its content.\par\medskip
\ifdefined\hypertarget\hypertarget{sourcepointer-3}{}\fi
\noindent\textit{Source reference.} Original lemma 016E: sources/stacks-hypercovering/pinned-original-source.tex, source lines 998--998. Cited by the retained statement or proof but not part of the collected spans; supplied only as a navigation pointer. The argument does not depend on its content.\par\medskip
\ifdefined\hypertarget\hypertarget{sourcepointer-4}{}\fi
\noindent\textit{Source reference.} Original lemma 018I: sources/stacks-hypercovering/pinned-original-source.tex, source lines 998--998. Cited by the retained statement or proof but not part of the collected spans; supplied only as a navigation pointer. The argument does not depend on its content.\par\medskip
\ifdefined\hypertarget\hypertarget{sourcepointer-5}{}\fi
\noindent\textit{Source reference.} Original lemma 018J: sources/stacks-hypercovering/pinned-original-source.tex, source lines 998--998. Cited by the retained statement or proof but not part of the collected spans; supplied only as a navigation pointer. The argument does not depend on its content.\par\medskip
\begin{lemma}
\label{lemma-acyclic-hypercover-sheaves}
Let $\mathcal{C}$ be a site.
Let $K$ be a simplicial presheaf.
Let $\mathcal{G}$ be a presheaf.
Let $K \to \mathcal{G}$ be an augmentation of $K$
towards $\mathcal{G}$. Assume that
\begin{enumerate}
\item The morphism of presheaves $K_0 \to \mathcal{G}$ becomes
a surjection after sheafification.
\item The morphism
$$
(d^1_0, d^1_1) :
K_1
\longrightarrow
K_0 \times_\mathcal{G} K_0
$$
becomes a surjection after sheafification.
\item For every $n \geq 1$ the morphism
$$
K_{n + 1} \longrightarrow (\text{cosk}_n \text{sk}_n K)_{n + 1}
$$
turns into a surjection after sheafification.
\end{enumerate}
Then $H_i(K) = 0$ for $i > 0$ and
$H_0(K) = \mathbf{Z}_\mathcal{G}^\#$.
\end{lemma}

\begin{proof}
Denote $K^n = \text{cosk}_n \text{sk}_n K$ for $n \geq 1$.
Define $K^0$ as the simplicial object with terms
$(K^0)_n$ equal to the $(n + 1)$-fold fibred product
$K_0 \times_\mathcal{G} \ldots \times_\mathcal{G} K_0$,
see Simplicial,
Example \ref{simplicial-example-fibre-products-simplicial-object}.
We have morphisms
$$
K \longrightarrow \ldots \to K^n \to K^{n - 1} \to \ldots \to K^1 \to K^0.
$$
The morphisms $K \to K^i$, $K^j \to K^i$ for $j \geq i \geq 1$ come
from the universal properties of the $\text{cosk}_n$ functors.
The morphism $K^1 \to K^0$ is the canonical morphism
from
Simplicial, Remark \ref{simplicial-remark-augmentation}.
We also recall that $K^0 \to \text{cosk}_1 \text{sk}_1 K^0$
is an isomorphism, see
Simplicial, Lemma \ref{simplicial-lemma-cosk-minus-one}.

\medskip\noindent
By Lemma \ref{lemma-compare-cosk0} we see that
$H_i(K^0) = 0$ for $i > 0$ and $H_0(K^0) = \mathbf{Z}_\mathcal{G}^\#$.

\medskip\noindent
Pick $n \geq 1$. Consider the morphism $K^n \to K^{n - 1}$.
It is an isomorphism on terms of degree $< n$.
Note that $K^n \to \text{cosk}_n \text{sk}_n K^n$ and
$K^{n - 1} \to \text{cosk}_n \text{sk}_n K^{n - 1}$
are isomorphisms. Note that $(K^n)_n = K_n$ and
that $(K^{n - 1})_n = (\text{cosk}_{n - 1} \text{sk}_{n - 1} K)_n$.
Hence by assumption, we have that $(K^n)_n \to (K^{n - 1})_n$
is a morphism of presheaves which becomes surjective after
sheafification. By Lemma \ref{lemma-acyclicity} we conclude that
$H_i(K^n) = H_i(K^{n - 1})$.
Combined with the above this proves the lemma.
\end{proof}

\end{document}
